\section{Discussion and Limitations}
\label{sec:discussion}

\paragraph{What verification guarantees.}
On PVS, RoadScribe's verifier and template fallback guarantee that every narrative agrees with its ledger.
The world-level results show that this is weaker than it sounds.
Uncalibrated, the ledger's road values carry an 18\% vehicle signature and a speed dependence that a single trip cannot remove.
A perfectly faithful narrator therefore contradicts another vehicle's narrator about the same street at most places.
The input-level results show the converse: a verifier can make a model's failures safe without making the model useful.
Both problems are visible before deployment, one with paired recordings and the other with the training set.

\paragraph{What the calibration needs.}
For agreement between vehicles, one gain per vehicle was enough. It could be estimated from about a kilometre of road shared with vehicles already calibrated, a condition that fleets and repeated commutes meet.
A vehicle without such data should be described only by ride statements, which are honest about depending on the vehicle.
The within-place elasticity ($\hat\gamma=0.86$ on all routes, interval [0.68, 1.02]) is higher than the 0.5 implied by the nominal road spectrum of ISO~8608.
Under the same model it corresponds to a waviness of about $w=2.7$.
The difference may also reflect nonlinear suspension behaviour or a speed-dependent driving style, and three vehicles cannot tell these apart.
Class agreement is insensitive to this choice (gain with $\gamma=0.5$: $\kappa=0.95$). Claims about the contribution of speed are not.

\paragraph{What the classes mean.}
Road classes are speed-normalized comfort classes of a dashboard sensor, not a pavement condition index.
Comfort at the seat differs from vibration at the dashboard.
Agreement between vehicles shows that a class is reproducible, not that it is correct.
The accuracy against PVS quality labels on held-out routes and vehicles (98.3\%, balanced 0.82) is evidence that it also tracks the labelled condition. These labels have three levels and are dominated by ``good'' (76\%).

\paragraph{Limitations.}
\begin{itemize}
\item \emph{Data.} PVS has three vehicle--driver pairs and one pass per vehicle and route. Vehicle and driver are therefore confounded, and the confidence intervals do not cover the sampling of vehicles.
\item \emph{Sensors.} The sensors are fixed modules, not phones, and the size of the vehicle signature depends on the mount (2.9--18.2\%).
\item \emph{Training data.} The language models were fine-tuned for one epoch on synthetic, template-style narratives.
\item \emph{Not evaluated.} We did not evaluate question answering, unseen audiences, reader preferences, or on-device latency and energy.
\item \emph{Event claims} are confirmed by both other vehicles only about half of the time.
\end{itemize}

\section{Conclusion}
\label{sec:conclusion}
A narrator that checks every sentence against its facts is only as trustworthy as those facts are identifiable, from the world and from the narrator's input.
For vehicle vibration, both levels can be tested before deployment.
On PVS, both tests flagged claims that a verifier had passed or would pass.
A gain per vehicle estimated on shared roads made road claims agree across vehicles nearly as well as two sensors in one car, and the language model's errors followed, type by type, the values its training input lacked.

\paragraph{Data and code availability.}
PVS is public~\cite{menegazzo2021pvs}.
The code, pre-registrations, analysis scripts and result files are available at \url{https://github.com/Nishkarsh5108/roadScribe}, including the decision log, the scripts behind every number in this paper and the inputs of the fine-tuning runs.

\paragraph{Disclosure of Interests.}
The authors have no competing interests to declare that are relevant to the content of this article.
